\subsection{对数映射}

引理 1. 假设 \(p \neq 0\)\(G\)\(G_1\)\(G\)。设 G 是一个李群，G＿1 是 G 的一个同构于标准群的开子群，且 \(x \in G\)。下列条件等价：

\begin{enumerate}
\def\labelenumi{(\roman{enumi})}
\item
  x 的某个幂属于 \(x\)\(\mathbf{G}_1\)；
\item
  存在一个严格递增的整数序列 \((n_{i})\)，使得当 \(x^{n_{i}}\)\(e\) 趋于 \(i\) 时，\(+\infty\) 趋于 e。
  (ii) \(\Rightarrow\) (i)：显然。
  (i) \(\Rightarrow\) (ii)：假设 \(y = x^{m} \in G_{1}\)。由第 5 号的命题 9 (i)，当 n 趋于 \(y^{p^n}\) 时，\(e\)\(n\)\(+\infty\) 趋于 e，换言之，当 n 趋于 \(x^{mp^n}\) 时，\(e\)\(n\)\(+\infty\) 趋于 e。
\end{enumerate}

命题 10. 假设 \(p \neq 0\)\(G\)\(G_f\)。设 G 是有限维李群。令 G＿f 为满足存在严格递增的整数序列 \(x \in G\) 使得当 \((n_i)\) 趋于 \(x^{n_i}\)\(e\) 时 \(i\) 趋于 e 的 \(+\infty\) 的集合。

\begin{enumerate}
\def\labelenumi{(\roman{enumi})}
\item
  \(G_{f}\) 在 G 中是开的。
\item
  存在唯一一个从 \(\psi\) 到 L(G) 的映射 \(G_{f}\)，具有如下性质：
\end{enumerate}

\begin{enumerate}
\def\labelenumi{(\alph{enumi})}
\item
  \(\psi (x^n) = n\psi (x)\) 对所有 \(x\in \mathbf{G}_f\) 和所有 \(n\in \mathbf{Z}\) 成立；
\item
  存在 e 在 \(\mathbf{V}\) 中的一个开邻域 \(e\)\(\mathbf{G}\)，使得 \(\psi |\mathbf{V}\) 是一个单射指数映射的逆映射。
\end{enumerate}

\begin{enumerate}
\def\labelenumi{(\roman{enumi})}
\setcounter{enumi}{2}
\tightlist
\item
  映射 \(\psi\) 是解析的。
\end{enumerate}

存在 G 的一个开子群同\(\mathbf{G}\)构于标准群（第 3 号，定理 4）。断言 (i) 随即由引理 1 得到。

令 U 是 L(G) 的一个开子群，\(\phi\colon\mathrm{U}\to\phi(\mathrm{U})\) 是 G 的一个指数映射，具有第 2 号命题 3 的性质。可以假设 U 足够小，使得 \(\phi(\mathrm{U})\subset\mathrm{G}_{f}\)。令 \(x\in\mathrm{G}_{f}\)。存在 \(m\in\mathbf{Z}-\{0\}\)，使得 \(x^{m}\in\phi(\mathrm{U})\)。元素 \(\frac{1}{m}\phi^{-1}(x^{m})\)\(m\) 不依赖于 m 的选择。事实上，若 \(m'\in\mathbf{Z}\) 使得 \(x^{m'}\in\phi(\mathrm{U})\)，则 \(x^{mm'}\in\phi(\mathrm{U})\)，并且

\[
m ^ {\prime} \phi^ {- 1} (x ^ {m}) = \phi^ {- 1} (x ^ {m m ^ {\prime}}) = m \phi^ {- 1} (x ^ {m ^ {\prime}}),
\]

令 \(\psi(x) = \frac{1}{m}\phi^{-1}(x^m)\)。则 \(\psi|\phi(U) = \phi^{-1}\)。另一方面，若 \(n \in \mathbf{Z}\)，则

\[
\psi (x ^ {n}) = \frac {1}{m} \phi^ {- 1} (x ^ {n m}) = \frac {n}{m} \phi^ {- 1} (x ^ {m}) = n \psi (x).
\]

因此 \(\psi\)\(x\) 具有命题的性质 (a) 和 (b)。在 x 的某个邻域内，\(\psi\) 由映射 \(x \mapsto x^m\)、\(y \mapsto \phi^{-1}(y)\) 和 \(z \mapsto \frac{1}{m} z\) 复合而成；故 \(\psi\) 在 \(G_f\) 上是解析的。

最后，设 \(\psi'\) 是从 \(G_f\) 到 L(G) 的映射，\(e\) 是 \(G_f\) 中 e 的一个邻域，满足对于 \(\psi'(x^n) = n\psi'(x)\) 和 \(x \in G_f\) 有 \(n \in \mathbf{Z}\)，且 \(\psi'|V'\) 是一个单射指数映射的逆映射。于是 \(\psi\) 与 \(\psi'\) 在 e 的某个邻域 W 上一致。若 \(e\)\(x \in G_f\)，则存在 \(n \in \mathbf{Z}\) 使得 \(x^n \in W\)。于是

\[
n \psi^ {\prime} (x) = \psi^ {\prime} (x ^ {n}) = \psi (x ^ {n}) = n \psi (x)
\]

从而 \(\psi = \psi'\)。

定义 2. 命题 10 的映射 \(\psi\) 称为 G 的对数映射，记作 \(\log_{\mathbb{G}}\)，或简记为 log。

命题 11. 假设 \(p \neq 0\)。设 x、y 是 G＿f 中的两个可交换元素。则 xy \(x, y\)\(G_f\)\(xy \in G_f\) 且 \(\log(xy) = \log x + \log y\)。

xy \(xy \in G_f\)\(U\)\(L(G)\)\(\phi: U \to \phi(U)\)\(G\)\(U\) 这一事实由引理 1 得到。令 U 是加法群 L(G) 的一个开子群，\(\log(\phi(U))\) 是 G 的一个指数映射，具有第 2 号命题 3 的性质；可以假设 U 足够小，使得 \(\phi\) 是 \(n \in \mathbf{Z} - \{0\}\) 的逆映射。选取适当的 \(x^n \in \phi(U)\)，使得 \(y^n \in \phi(U)\)、\(u = \log x^n\)。令 \(v = \log y^n\)、\(x^n = \phi(u)\)，于是 \(y^n = \phi(v)\)、\([u,v] = 0\)。由公式 (2)，\(\phi(\lambda(u + v)) = \phi(\lambda u)\phi(\lambda v)\)。于是由 Hausdorff 公式，当 \(|\lambda|\)\(i\) 充分小时，；因而对于每个充分大的整数 i，

\[
\phi (p ^ {i} (u + v)) = \phi (p ^ {i} u) \phi (p ^ {i} v)
\]

即

\[
p ^ {i} (\log x ^ {n} + \log y ^ {n}) = \log (x ^ {n p ^ {i}} y ^ {n p ^ {i}})
\]

或

\[
n p ^ {i} (\log x + \log y) = n p ^ {i} \log (x y).
\]

命题 12. 假设 \(p \neq 0\)。设 \(x \in G_f\)。下列条件等价：

\begin{enumerate}
\def\labelenumi{(\roman{enumi})}
\item
  \(\log x = 0;\)
\item
  x 在 G 中是有限阶的。
\end{enumerate}

若存在整数 \(n > 0\) 使得 \(x^n = e\)，则

\[
n \log x = \log x ^ {n} = 0,
\]

若 \(\log x = 0\)\(\log x = 0\)，令 \(\mathrm{V}\)\(e\) 为 \(\mathbf{G}_f\) 中 e 的一个邻域，使得 \(\log |\mathrm{V}\) 是一个单射指数映射的逆映射。存在整数 \(n > 0\) 使得 \(x^n \in \mathrm{V}\)；等式 \(\log x^n = 0\) 蕴含 \(x^n = e\)。

命题 13. 假设 \(p \neq 0\)\(G\)。若 G 是紧群或标准群，则 \(G_f = G\)。

若 G 是标准群，只需应用引理 1。假设 G 是紧群。设 \(x \in G\)，V 是 G 中 e 的一个邻域。令 y 为序列 \(e\)\(y\)\((x^n)_{n \geqslant 0}\) 的一个聚点。对于每个 \(n > 0\)，存在两个整数 \(n_1, n_2\)，使得 \(n_1 \geqslant 2n_2 \geqslant n\)，且 \(x^{n_1} \in y\mathrm{V}, x^{n_2} \in y\mathrm{V}\)\(x^{n_1 - n_2} \in V^{-1}\mathrm{V}\)、，从而 ，并且 \(n_1 - n_2 \geqslant n\)。因此 \(x \in G_f\)。

推论。假设 K 局部紧。则 \(G_{f}\) 是 G 的紧子群的并。

令 \(x \in G\)\(x\)\(G\)。若 x 属于 G 的某个紧子群，则 \(x \in G_f\)（命题 13）。假设 \(x \in G_f\)\(K\)。由于 K 局部紧，G 存在一个紧的开子群 \(G_1\)\(G\)。于是存在整数 \(m > 0\)，使得 \(x^m \in G_1\)。由 \(G_2\) 生成的闭子群 \(x^m\) 包含于 \(G_1\)\(x\)，因而是紧的。于是 x 与 \(G_2\) 的元素交换，从而 \(G_2 \cup xG_2 \cup \cdots \cup x^{m-1}G_2\)\(G\)\(x\) 是 G 的包含 x 的紧子群。

例。假设 K 局部紧。令 U 为 A 的可逆元素的集合；它是李群 K* 的紧开子群。由命题 13，\(U \subset (\mathbf{K}^{*})_{f}\)；另一方面，若 \(x \in K^{*}\) 且 \(x \notin U\)，则当 n 趋于 \(x^{n}\) 时 \(+\infty\) 趋于 0，或者当 n 趋于 \(x^{n}\) 时 \(-\infty\) 趋于 0；因此 \(U = (\mathbf{K}^{*})_{f}\)。函数 \(\log_{K^{*}}\) 在 \(U_{f}\) 上有定义且解析，取值于 \(L(\mathbf{K}^{*}) = \mathbf{K}\)，并满足对 U 中所有 x、y，\(\log_{\mathbf{K}^{*}}(xy) = \log_{\mathbf{K}^{*}}(x) + \log_{\mathbf{K}^{*}}(y)\)；U 中满足 \(\log_{\mathbf{K}^{*}}(x) = 0\) 的元素 x 正是 K 的单位根。

我们再次采用第 3、4、5 号的记号。

命题 14. 假设 \(p \neq 0\)，且 \(v\) 选取为满足 \(v(p) = 1\)\(G\)\(E(X)\)\(L(X)\)\(G\)\(G\)。设 G 是标准群，E(X)（分别地，L(X)）是 G 的指数函数（分别地，对数函数）在 0 点的整级数展开。

\begin{enumerate}
\def\labelenumi{(\roman{enumi})}
\item
  E 的绝对收敛域（可微与解析流形，R，4.1.3）包含集合 \(\Delta\)，其中 \(x \in G\) 满足 \(\omega(x) > \frac{1}{p - 1}\)。令 \(\mathbf{E}'\) 表示由该级数在 \(\Delta\) 上定义的映射。则 \(\mathbf{E}'\) 是 G 的指数映射，且是从流形 \(\Delta\) 到自身的同构。
\item
  L 的绝对收敛域包含 G。令 L' 表示由该级数在 G 上定义的映射。则 L' 是 G 的对数映射，且 L' 在 \(\Delta\) 上的限制是 E' 的逆映射。
\item
  映射 \(\mathbf{E}'\) 在带 Hausdorff 法则的 \(\Delta\) 上是同构，并且其像为 \(\Delta\) 的子群 \(\mathbf{G}\)。
\end{enumerate}

采用 § 5，第 3、4 号的记号，有 \(E = \sum_{m \geqslant 1} \frac{\psi_{m,m}}{m!}\)（§ 5，第 4 号，命题 3）。由于系数 \(c_{\alpha\beta\gamma}\) 属于 A，\(\|\psi_{m,m}\| \leqslant 1\)（可微与解析流形，R，附录）。假定 \(K^{r}\) 采用范数

\[
\left\| \left(\lambda_ {1}, \dots , \lambda_ {r}\right) \right\| = \sup \left(\left| \lambda_ {1} \right|, \dots , \left| \lambda_ {r} \right|\right).
\]

由第二章，§ 8，第 1 号，引理 1，有 \(v(m!) \leqslant \frac{m - 1}{p - 1}\)。若 \(\omega(x) > \frac{1}{p - 1}\)，则可见当 m 趋于无穷时，\(m\omega(x) - v(m!)\) 趋于 \(+\infty\)\(m\)，从而

\[
\left\| \frac {\psi_ {m , m}}{m !} \right\| \| x \| ^ {m} \leqslant \frac {1}{| m ! |} \| x \| ^ {m}
\]

当 m 趋于  时，它趋于 0，并且

\[
\omega \left(\frac {\psi_ {m , m} (x)}{m !}\right) > \frac {m}{p - 1} - \frac {m - 1}{p - 1} = \frac {1}{p - 1} \quad \text { for } m \geqslant 1.
\]

因此 \(\Delta\) 包含于 E 的绝对收敛域，且 \(\mathrm{E}'(\Delta) \subset \Delta\)。显然，\(\mathrm{E}'\) 是一个指数映射。

若 \(L_{m}\)\(L\)\(m\) 表示 L 的 m 次齐次分量，则 §5，第 4 号，命题 3 说明，\(L_{m}\) 的每个系数都具有如下形式

\[
a _ {1} + \frac {1}{2} a _ {2} + \dots + \frac {1}{m} a _ {m}
\]

其中 \(a_1, a_2, \ldots, a_m\) 属于 A；但是

\[
\inf \left(v (1), v \left(\frac {1}{2}\right), \dots , v \left(\frac {1}{m}\right)\right) = o (\log m) \quad \text { as   } m \text {   tends   to   } + \infty
\]

且

\[
\inf \left(v (1), v \left(\frac {1}{2}\right), \dots , v \left(\frac {1}{m}\right)\right) \geqslant v \left(\frac {1}{m !}\right) \geqslant - \frac {m - 1}{p - 1}.
\]

因此，若 \(\omega(x) > 0\)，则 \(\|L_m\|. \|x\|^m\)\(m\) 当 m 趋于  \(+\infty\) 时趋于 0，所以 G 包含于 L 的绝对收敛域。另一方面，若 \(\omega(x) > \frac{m}{p-1}\)，则当 \(\omega(L_m(x)) > \frac{m}{p-1} - \frac{m-1}{p-1} = \frac{1}{p-1}\) 时，\(m \geqslant 1\)，从而 \(L'(\Delta) \subset \Delta\)。

由于形式幂级数 \(\mathrm{L}(\mathrm{E}(\mathrm{X}))\) 和 \(\mathrm{E}(\mathrm{L}(\mathrm{X}))\) 都等于 \(\mathbf{X}\)，可微与解析流形，\(\mathbf{R}\)，4.1.5 说明

\[
\mathrm{L} ^ {\prime} (\mathrm{E} ^ {\prime} (x)) = \mathrm{E} ^ {\prime} (\mathrm{L} ^ {\prime} (x)) = x
\]

对 \(x \in \Delta\)。因此 \(E'\) 是从流形 \(\Delta\) 到自身的同构，其逆同构是 \(L'\) 在 \(\Delta\) 上的限制。

\(\mathrm{L}(\mathrm{X}^{[n]}) = n\tilde{\mathrm{L}}(\mathrm{X})\) 对于整数 \(n\) 大于 \(> 0\) 成立（参见 § 5，第 4 号）。由于 \(G\) 以及 \(L\) 都包含于 \(X^{[n]}\) 的绝对收敛域，因而对所有 \(L'(x^n) = nL'(x)\)，有 \(x \in G\)。关系 \(L'| \Delta = E'^{-1}\) 蕴含，当 \(L'(x^n) = \log x^n\) 足够大时 \(n\)。所以 \(L'(x) = \log x\)。至此已经证明 (i) 和 (ii)。

令 \(H = \sum_{r,s \geq 0} H_{r,s}\) 为 Hausdorff 形式幂级数，h 为相对于 L(G) 的 Hausdorff 函数。\(\tilde{H}\) 的绝对收敛域包含 \(\Delta \times \Delta\)\(h\)，且 h 定义在 \(\Delta \times \Delta\) 上（第二章，§ 8，命题 2）。于是

\[
\mathrm{E} ^ {\prime} (x) \mathrm{E} ^ {\prime} (y) = \mathrm{E} ^ {\prime} (h (x, y))
\]

当 \(x, y\) 充分接近 0 时（§ 4，定理 4 (v)）。因此，按照第 3 号定义 1 的记号，形式幂级数 \(\mathbf{F}(\mathrm{E}(\mathbf{X}), \mathrm{E}(\mathbf{Y}))\) 和 \(\mathrm{E}(\mathrm{H}(\mathbf{X}, \mathbf{Y}))\) 相等。设 \(x, y\) 是 \(\Delta\) 的元素。于是

\[
\begin{array}{l} \sup _ {m} \left\| \frac {\psi_ {m , m}}{m !} \right\| (\sup \| x \|, \| y \|) ^ {m} <   1 \\ \sup _ {r, s} \| H _ {r, s} \| \| x \| ^ {r} \| y \| ^ {s} <   | p | ^ {1 / (p - 1)} \end{array}
\]

由第二章，§8，公式 (14)，以及可微与解析流形，R，4.1.5，\(\mathrm{E}'(x)\mathrm{E}'(y)\) 是将 \(x\) 代入 X、将 \(y\) 代入 Y 后得到的，而

\[
\mathrm{F} (\mathrm{E} (\mathrm{X}), \mathrm{E} (\mathrm{Y}))
\]

且 \(\mathrm{E}'(h(x,y))\) 是将 \(x\) 代入 \(\mathbf{X}\)、将 \(y\) 代入 \(\mathbf{Y}\) 后在 \(\mathrm{E}(\mathrm{H}(\mathbf{X},\mathbf{Y}))\) 中得到的。因此 \(\mathrm{E}'(x)\mathrm{E}'(y) = \mathrm{E}'(h(x,y))\)。

